\documentclass[12pt,a4paper]{article}

% ------------------------
% Paquetes
% ------------------------
\usepackage[spanish]{babel}
\usepackage[utf8]{inputenc}
\usepackage[T1]{fontenc}
\usepackage{lmodern}
\usepackage{geometry}
\usepackage{enumitem}
\usepackage{booktabs}
\usepackage{hyperref}
\geometry{margin=2.5cm}

% ------------------------
% Datos del documento
% ------------------------
\title{\textbf{Concurrent Computing}\\
\large Licenciatura en Ciencia de Datos}
\author{}
\date{}

\begin{document}
\maketitle

First of all, we would like to know wat even concurrence means. Other keyword is process, threads, we'll see the difference between them. We'll know hiperthreading in a processor.
There are lots of branches, like sequential (our generally first approach), concurrent, parallel and even distributed.

For the third (last) moduli, we'll use a framework based in python named Django.
The evaluation is as follows

Activities
Exams for each moduli
Expositions
Practices
Projects
Participations

The rounding is n+0.6=n+1

Practices:

In the email we need to have in the title the practice number and the team number, for example: P1E1

They are contempled in PDF, with commands, code, images and the printing of screens that reflexes the correctness of the practice.
The practice name is similar, practice number + team number.

The archive has to be attached instead of being part of the message. In the case we need to deliver more than one archive then we have to make a zip.

Software:
The OS doesn't matter, it could ve windows, linux or macOS; we'll use python, MPI (message passing interface) and virtual box.

chavesti1@gmail.com


\end{document}

